% Internal and entrusted manufacturing technologies.
\section{Internal and Entrusted Manufacturing Technologies}
\label{sec:p03-technologies}

This section defines the manufacturing primitives that later enter route
values. It does not construct those values, choose a route, or clear the
qualified manufacturing-service market. All technology primitives are
invariant to the MAH regime. The only direct institutional wedge remains
\(\tau_E(M)\).

\subsection{Internal manufacturing}
\label{subsec:p03-internal}

For a project with manufacturing requirement \(m>0\) and a developer with
internal manufacturing capability \(k_i>0\), let
\begin{equation}
  c_I(m,k_i)>0,
  \qquad
  c_{I,m}(m,k_i)>0,
  \qquad
  c_{I,k}(m,k_i)<0
  .
  \label{eq:p03-internal-cost}
\end{equation}
The function \(c_I\) is technological marginal manufacturing cost, measured
in \(\mathsf C/\mathsf Y\). A more demanding project raises this cost, while
greater internal capability lowers it on the baseline domain.

Internal production also requires a finite project-level
production-readiness/setup cost \(F_I(m,k_i)\), measured in \(\mathsf C\):
\begin{equation}
  0\leq F_I(m,k_i)<+\infty,
  \qquad
  F_{I,m}(m,k_i)>0,
  \qquad
  F_{I,k}(m,k_i)<0.
  \label{eq:p03-internal-setup}
\end{equation}
Low internal capability therefore makes internalization expensive but does
not make it technologically impossible by convention.
The setup term is the project-level present value of building, acquiring,
upgrading, validating, or adapting internal production capability, including
internal quality-system preparation and project-specific fixed
organizational costs. It may become very large as \(k_i\) approaches zero.
The optional limiting condition
\(\lim_{k\downarrow0}F_I(m,k)=+\infty\) is not a discrete feasibility cutoff.

\subsection{Qualified external production}
\label{subsec:p03-external}

Under entrusted manufacturing, a qualified external producer uses the
technological marginal-cost kernel
\begin{equation}
  c_E(m)>0,
  \label{eq:p03-external-cost}
\end{equation}
also measured in \(\mathsf C/\mathsf Y\). This kernel does not depend on the
developer's internal capability \(k_i\): manufacturing is performed by the
qualified external producer.

An entrusted project requires
\begin{equation}
  b(m)>0,
  \qquad b'(m)>0,
  \label{eq:p03-capacity-requirement}
\end{equation}
units of qualified manufacturing-service capacity. The object \(b(m)\) is a
physical capacity requirement in \(\mathsf B\) per project. It is not a route
share, a supply curve, or an exogenous measure of scarcity relief.

The real technology-transfer, validation, and production-readiness cost is
\begin{equation}
  F_E(m)\geq0,
  \label{eq:p03-external-fixed-cost}
\end{equation}
measured in \(\mathsf C\) per project. In addition, the developer bears a
residual holder-side responsibility and coordination burden
\begin{equation}
  \mu_E\geq0.
  \label{eq:p03-holder-burden}
\end{equation}
The burden \(\mu_E\) is not eliminated by MAH. Entrusted production avoids
full proprietary manufacturing capability but still requires technology
transfer, validation, contracted qualified capacity, and holder-side quality
oversight.

The accounting roles are distinct. The technological marginal cost \(c_E(m)\)
is the input to the operating-value kernel derived above. The separate
monetary capacity payment \(p_m b(m)\), the readiness cost \(F_E(m)\), the
holder burden \(\mu_E\), and the institutional barrier \(\tau_E(M)\) are not
embedded in \(c_E(m)\). Each can therefore enter a later route value at most
once.

\subsection{Organizational and institutional distinction}
\label{subsec:p03-organization}

Routes \(I\) and \(E\) are technologically and organizationally distinct:
\begin{itemize}
  \item under \(I\), the developer remains the holder and manufactures using
  its own capability \(k_i\);
  \item under \(E\), the developer remains the authorization holder, retains
  holder responsibility, and a qualified external producer manufactures;
  \item route \(E\) is not the transfer or out-licensing route \(T\), under
  which the retained holder--producer organization is not used.
\end{itemize}

The only direct policy object is
\begin{equation}
  M\in\{0,1\},
  \qquad
  \tau_E(0)=+\infty,
  \qquad
  \tau_E(1)=\bar\tau_E<+\infty.
  \label{eq:p03-policy-invariance}
\end{equation}
The regime does not change \(c_I\), \(F_I\), \(c_E\),
\(b\), \(F_E\), \(\mu_E\), \(a_i\), \(k_i\), \(q\), \(m\), \(s(q)\), or
\(s_g(q)\). In particular, this technology block does not posit a fall in
\(p_m^*\) or an outward shift of CMO supply. Price and capacity are equilibrium
objects determined in the qualified manufacturing-service market below.

\subsection{Boundary checks}
\label{subsec:p03-boundaries}

Increasing \(k_i\) lowers internal marginal and setup costs, while increasing
\(m\) raises both. A low-\(k_i\) developer can still internalize by paying the
high real cost. The entrusted technology does not use \(k_i\), while a larger
\(m\) raises its qualified-capacity requirement \(b(m)\). None of these
statements is a route-choice result: route values, outside options, and
deterministic sorting are defined in the next section.
